% PAGES: 277-277
% Closes the \begin{itemize} opened at the end of sec09_c1_10.tex (PDF page
% 276) with the second bullet below (weight w_{min,cash}, eq. 9.34).

\item weight $w_{min,cash}$ of the cash position given by
\begin{equation}
w_{min,cash} \;=\; 1-\bfone^tw_{min} \;=\;
1-(\mu_P-r_f)\frac{\bfone^t\Sigma_R^{-1}\bar\mu}{\bar\mu^t\Sigma_R^{-1}\bar\mu}.
\end{equation}
\end{itemize}

The minimum standard deviation of the return of a portfolio with given expected return
$\mu_P$ is
\begin{align}
\sigma_{min} &= \sqrt{w_{min}^t\Sigma_Rw_{min}} \;=\;
\frac{\mu_P-r_f}{\bar\mu^t\Sigma_R^{-1}\bar\mu}
\sqrt{(\Sigma_R^{-1}\bar\mu)^t\Sigma_R(\Sigma_R^{-1}\bar\mu)} \notag\\
&= \frac{\mu_P-r_f}{\bar\mu^t\Sigma_R^{-1}\bar\mu}
\sqrt{\bar\mu^t(\Sigma_R^{-1})^t(\Sigma_R\Sigma_R^{-1})\bar\mu} \;=\;
\frac{\mu_P-r_f}{\bar\mu^t\Sigma_R^{-1}\bar\mu}\sqrt{\bar\mu^t\Sigma_R^{-1}\bar\mu} \notag\\
&= \frac{\mu_P-r_f}{\sqrt{\bar\mu^t\Sigma_R^{-1}\bar\mu}};
\end{align}
here, we used the facts that $(\Sigma_R^{-1})^t = \Sigma_R^{-1}$, since $\Sigma_R$ is a
symmetric matrix; see Lemma 1.9. Note that $\bar\mu^t\Sigma_R^{-1}\bar\mu > 0$ since
$\Sigma_R$ is a symmetric positive definite matrix, and therefore $\Sigma_R^{-1}$ is
also symmetric positive definite; cf. Lemma 5.4.

The pseudocode for the asset allocation for the minimum variance portfolio computed
using the formulas (9.33--9.35) can be found in Table 9.1.

\par\medskip\noindent\textit{Example:}\ Find the \$100 million minimum variance
portfolio with expected return equal to 2.25\% and invested in cash and four assets
with expected returns equal to 2\%, 1.75\%, 2.5\%, and 1.5\%, and with the following
covariance matrix of their returns:
\[
\begin{pmatrix}
0.09 & 0.01 & 0.03 & -0.015 \\
0.01 & 0.0625 & -0.02 & -0.01 \\
0.03 & -0.02 & 0.1225 & 0.02 \\
-0.015 & -0.01 & 0.02 & 0.0576
\end{pmatrix}.
\]

Assume that the risk free rate is 1\%.

\begin{answerx}
% Left OPEN at the end of this file: the derivation continues onto PDF page
% 278 (folio 262), which closes it with "...respectively." followed by the
% printed square. sec09_c1_12.tex must NOT reopen \begin{answerx}.
Let $\mu_1 = 0.02$, $\mu_2 = 0.0175$, $\mu_3 = 0.025$, $\mu_4 = 0.015$, and $r_f =
0.01$. We look for the asset and cash allocation corresponding to the minimum variance
portfolio with expected return $\mu_P = 0.0225$.

Note that $\bar\mu = \mu - r_f\bfone$ and $\Sigma_R$ are given by
\[
\bar\mu \;=\; \begin{pmatrix} 0.01 \\ 0.0075 \\ 0.015 \\ 0.005 \end{pmatrix}; \quad
\Sigma_R \;=\; \begin{pmatrix}
0.09 & 0.01 & 0.03 & -0.015 \\
0.01 & 0.0625 & -0.02 & -0.01 \\
0.03 & -0.02 & 0.1225 & 0.02 \\
-0.015 & -0.01 & 0.02 & 0.0576
\end{pmatrix}.
\]

The vector $\Sigma_R^{-1}\bar\mu$ is computed by using the Cholesky decomposition to
solve the linear system corresponding to the matrix $\Sigma_R$ and to the right hand
side $\bar\mu$, i.e., $\Sigma_R^{-1}\bar\mu =
\text{linear\_solve\_cholesky}(\Sigma_R,\bar\mu)$; see Table 6.2 and the pseudocode
from Table 9.1. From (9.33) and (9.34), we find that the asset weights vector
corresponding to the minimum variance portfolio is
\[
w_{min} \;=\; \frac{\mu_P-r_f}{\bar\mu^t\Sigma_R^{-1}\bar\mu}\Sigma_R^{-1}\bar\mu \;=\;
\begin{pmatrix} 0.2120 \\ 0.4888 \\ 0.3540 \\ 0.2808 \end{pmatrix},
\]
